All primary solves and all 168 closed-loop episodes completed without recorded solver failures. The maximum true-model bounded-noise successor and selected-set robust successor were both $0.9999999701$ or lower; all inputs satisfied their bounds to numerical tolerance. The 18,432 DCC evaluations satisfied the gap thresholds, sandwich inequalities, action bounds, and post-compaction support counts. Maximum lower-bound excess over the full-data objective was $4.1\times10^{-14}$, and maximum excess over the action-error bound was $1.7\times10^{-13}$. These residuals support numerical consistency, not exact floating-point certification.

\paragraph{Conservatism and retained information.}
At $\epsilon=10^{-3}$, mean excess cost was $6.94\times10^{-6}$ over all plants; the maximum was $9.73\times10^{-4}$. At $10^{-5}$, the mean was $8.99\times10^{-9}$ and maximum below $7\times10^{-6}$. Table~\ref{tab:main} and Figure~\ref{fig:performance} separate plants. DCC retained a mean of 44.0 scalar data inequalities at $10^{-3}$, versus 11,264 for full histories averaged over plants, plus the excluded fixed prior. The largest observed count was 152, below each plant's $2n(n+m)$ bound. DCC retained at most two witnesses for the two-input plants and one for the three-input dense plant.

Fixed-budget data subsets had larger average excess costs under uniform noise: $0.0255$ for random, $0.0189$ for log-determinant, and $0.0275$ for rolling-window selection. Under boundary noise all three matched the full-data filter up to numerical error. This negative control matters: the method is not uniformly more accurate than subset heuristics. Here, informative boundary samples made even small subsets effectively exact. DCC supplies a per-query certificate in either regime rather than assuming that informativeness.

\paragraph{Archive access and timing.}
At $10^{-3}$, queries accepted without archive calls ranged from $29.6\%$ for the nonnormal plant to $91.2\%$ for the dense plant. Mean total times were 22.5, 12.4, 55.6, and 15.1\,ms, versus 48.7, 35.4, 111.9, and 100.1\,ms for full data. DCC was generally slower than fixed-budget heuristics, except in the dense case. Its tradeoff is certified conservatism and lower retained information, not universal speed. These local wall times include archive work, omit real network latency, and are affected by concurrent computation. Appendix tables expose the local/archive split.

\paragraph{Scaling, reuse, and streaming.}
Scaling retained at most 15, 47, 139, and 256 inequalities for $n=2,4,8,12$, below bounds 16, 48, 160, and 336. At $n=12$, increasing the archive from 128 to 8,192 transitions increased mean full-data time from 32.7\,ms to 2.83\,s and DCC total time from 25.8\,ms to 0.793\,s. Archive computation still grows with history length (Figure~\ref{fig:cert}).

Resetting the cache before each ablation query used 8.33 archive LPs per query on average; reuse with compaction used 6.45. Disabling compaction reduced this to 2.47 but increased retained data from 29.4 to 44.5 inequalities on average over short 32-query sequences. Aggressive compaction can therefore sacrifice future reuse. An input-box-only lower bound, omitting model witnesses, gave a mean certificate gap of $0.474$, although the two-sided gap was below $10^{-4}$. Optimistic witness-only actions violated a full-model face in eight ablation queries (maximum $2.78\times10^{-5}$), demonstrating why a lower certificate alone is insufficient for safety.

Across 720 streaming queries, screening invalidated 41 cached row witnesses. Actions reconstructed from exported payloads alone agreed with archive-backed actions within $3\times10^{-14}$; all reconstructed gaps met $10^{-4}$ to tolerance. At 1,024 transitions, the median JSON payload/raw-float64-data ratio was $0.104$; short histories can have ratios above one (Figure~\ref{fig:ablation}). The stale-witness example produced a near-zero claimed gap but an actual gap of $0.06$ before screening.

\paragraph{Closed loop.}
All 16,128 physical control steps remained in the safe box, with maximum coordinate magnitude below $0.9999999701$. Mean cumulative intervention per 96-step episode was 139.713411 for full data and DCC at $10^{-5}$, and 139.713333 for DCC at $10^{-3}$. The latter can be smaller because policy trajectories differ; it is not a violation of common-query optimality. Prior-only control had mean intervention 139.760339. These small differences indicate that primary stress-query improvements do not translate into a large closed-loop intervention improvement in this deployment experiment.
